\begin{abstract}
For the closed disk $D_R\subset\R^2$, let
\[
W(R):=\sup\{\Vol(B):B\text{ is wrappable by }D_R\}
\]
in the $1$-Lipschitz wrapping formalization of Nandakumar.  Scaling reduces the problem to the single dimensionless constant $W(1)$.  A universal area bound and the Euclidean isoperimetric inequality give
\[
W(R)\le \frac{\pi}{6}R^3.
\]
We also isolate the rotational-radial subclass: solids of revolution generated by a simple north--south meridian of length at most $R$, wrapped by the direct radial map.  Allowing the meridian to be nonconvex does not improve this subclass.  Its sharp value is the classical Mylar value
\[
W_{\rm rev}^{\rm rad}(R)
=\frac{\pi}{12I^2}R^3
=0.152315527\ldots R^3,
\qquad
I=\int_0^1\frac{du}{\sqrt{1-u^4}}.
\]
We then give a $17$-vertex, $30$-face non-axisymmetric polyhedral wrapping with exact volume
\[
C R^3,
\qquad
C=0.17984349416291615\ldots,
\]
certified by two independent exact-arithmetic programs.  Consequently
\[
W_{\rm rev}^{\rm rad}(R)<C R^3\le W(R),
\]
so the unrestricted problem strictly exceeds the sharp rotational-radial optimum by about $18.07\%$.  In this sense the construction exhibits a symmetry-breaking phenomenon at the level of optimal values.  Nandakumar's qualitative Question~1 is much weaker: under the same formalization it has a short elementary affirmative answer by making an arbitrarily small inward rotational notch in a suitable oblate spheroid.  The main point here is therefore quantitative: to begin locating the extremal constant $W(1)$ and to show that sharp rotational symmetry is already suboptimal.
\end{abstract}

\section{Introduction}
Nandakumar recently formulated a volume-maximization problem for wrapping three-dimensional solids by a planar sheet that may be folded or wrinkled but not stretched or torn~\cite{Nandakumar2026}.  In his formalization, a planar region $D$ wraps a compact connected solid when its boundary surface is realized by a $1$-Lipschitz image of $D$, with the boundary of $D$ glued to itself in the image.  For the disk $D_R$ of radius $R$, define
\begin{equation}\label{eq:Wdef}
W(R):=\sup\{\Vol(B):B\text{ is wrappable by }D_R\}.
\end{equation}
Scaling source and target shows immediately that
\begin{equation}\label{eq:scaling}
W(R)=R^3W(1),
\end{equation}
so the unrestricted disk-wrapping problem is governed by a single dimensionless constant $W(1)$.

Nandakumar's Question~1 asks whether $D_\pi$ can wrap a nonconvex solid of volume greater than the unit ball $4\pi/3$~\cite[Question~1]{Nandakumar2026}.  Under the stated $1$-Lipschitz formalization, this qualitative threshold is not difficult: \cref{sec:soft} gives a short elementary construction by adding a tiny inward rotational notch to an oblate spheroid whose volume is already larger than $4\pi/3$ and whose north--south meridian length is strictly less than $\pi$.  Thus nonconvexity by itself is not the central quantitative obstacle.

The more informative question is how large $W(1)$ actually is.  There are two elementary reference points.  First, every wrappable boundary has area at most the area $\pi R^2$ of the source disk, and the Euclidean isoperimetric inequality therefore gives the universal upper bound
\begin{equation}\label{eq:universalupperintro}
W(R)\le \frac{\pi}{6}R^3.
\end{equation}
Second, the classical Mylar variational problem solves a natural rotational subclass exactly.  Let $W_{\rm rev}^{\rm rad}(R)$ denote the supremum of volumes of solids obtained by revolving a simple north--south meridian of length at most $R$, with the disk mapped to the resulting surface by the direct radial parametrization of \cref{lem:radial}.  In \cref{sec:rotational} we show, without any convexity assumption on the meridian, that
\begin{equation}\label{eq:rotintro}
W_{\rm rev}^{\rm rad}(R)
=\frac{\pi}{12I^2}R^3
=0.152315527\ldots R^3,
\qquad
I=\int_0^1\frac{du}{\sqrt{1-u^4}}.
\end{equation}
This is exactly the Mylar value.  The inequality is classical in substance; we include a short global proof because its role as a sharp barrier for disk wrapping is central here.

Our main construction crosses that barrier by a substantial amount.

\begin{theorem}[Certified separation from the rotational-radial optimum]\label{thm:main}
For every $R>0$,
\begin{equation}\label{eq:mainchain}
W_{\rm rev}^{\rm rad}(R)
=\frac{\pi}{12I^2}R^3
< C R^3
\le W(R)
\le \frac{\pi}{6}R^3,
\end{equation}
where
\begin{equation}\label{eq:Cexact}
C=
\frac{
359686988325832312742347877865122384142690317017149
}{
2000000000000000000000000000000000000000000000000000
}
=0.17984349416291615\ldots .
\end{equation}
More precisely, there exists a continuous $1$-Lipschitz map from $D_R$ onto an embedded nonconvex polyhedral $2$-sphere with $17$ vertices and $30$ triangular faces, bounding a solid of volume $C R^3$.
\end{theorem}

The strict middle inequality in \eqref{eq:mainchain} does not depend on decimal rounding.  Indeed,
\[
I>\int_0^{3/4}1\,du
+\int_{3/4}^1\frac{du}{2\sqrt{1-u}}
=\frac54,
\]
because $1-u^4=(1-u)(1+u+u^2+u^3)<4(1-u)$ for $u<1$.  Hence
\[
\frac{\pi}{12I^2}<\frac{4\pi}{75}<0.168<C.
\]

The ratio
\[
\frac{C}{\pi/(12I^2)}=1.1807\ldots
\]
shows that the certified construction exceeds the sharp rotational-radial optimum by about $18.07\%$.  The significance of the decimal $0.179843\ldots$ is therefore not that this particular number is expected to be optimal.  Rather, it lies strictly beyond the exact threshold $0.152315527\ldots$ for the entire rotational-radial class.  At the level of optimal values, rotational symmetry is genuinely suboptimal.

The witness was discovered by numerical search, but \cref{thm:main} depends only on the final finite exact certificate.  We give the source triangulation, target coordinates, and piecewise-affine map explicitly.  Two supplied implementations use different exact arithmetic and different triangle-intersection algorithms.  They independently certify strict nonexpansion on all $30$ source triangles, the topology of the quotient complex, all $435$ triangle-pair intersections, nonconvexity, and the exact enclosed volume.

The paper is organized around the extremal problem.  \Cref{sec:bounds} records scaling and the universal isoperimetric upper bound.  \Cref{sec:soft} disposes of the qualitative nonconvexity threshold in Question~1.  \Cref{sec:rotational} establishes the sharp rotational-radial benchmark.  The remaining sections specify and certify the non-axisymmetric construction that crosses this benchmark.

\section{The extremal quantity and a universal upper bound}\label{sec:bounds}
The scaling relation \eqref{eq:scaling} follows by dilating both source and target.  Thus it is enough to understand $W(1)$.  The following coarse upper bound places the problem in a finite interval.

\begin{proposition}[Universal isoperimetric upper bound]\label[proposition]{prop:upper}
For every $R>0$,
\[
W(R)\le \frac{\pi}{6}R^3.
\]
Equivalently, $W(1)\le\pi/6$.
\end{proposition}

\begin{proof}
Let $B$ be any solid wrappable by $D_R$, and let $S=\partial B$ be the closed surface appearing as the image of a $1$-Lipschitz wrapping map $F:D_R\to S$.  A $1$-Lipschitz map cannot increase $2$-dimensional Hausdorff measure, hence
\[
\mathcal H^2(S)\le \mathcal H^2(D_R)=\pi R^2.
\]
The Euclidean isoperimetric inequality in $\R^3$ gives
\[
36\pi\,\Vol(B)^2\le \mathcal H^2(S)^3.
\]
Combining the two inequalities yields
\[
\Vol(B)\le
\frac{(\pi R^2)^{3/2}}{6\sqrt\pi}
=\frac{\pi}{6}R^3.
\]
Taking the supremum over all wrappable $B$ proves the claim.
\end{proof}

The bound $\pi/6=0.5235987755\ldots$ is intentionally crude; it uses only the area resource of the source disk and ignores the much stronger metric restrictions imposed by being a $1$-Lipschitz quotient of a flat disk.  Narrowing the gap between this upper bound and the certified lower bound in \cref{thm:main} is a natural continuation of the problem.

\section{A soft affirmative answer to Nandakumar's Question~1}\label{sec:soft}
We first isolate a simple observation that applies to convex and nonconvex surfaces of revolution alike.

\begin{lemma}[Radial wrapping of a surface of revolution]\label[lemma]{lem:radial}
Let
\[
\gamma(s)=(\varrho(s),z(s)),\qquad 0\le s\le L,
\]
be a simple rectifiable arc in the half-plane $\varrho\ge0$, parametrized by arclength, with
\[
\varrho(0)=\varrho(L)=0,
\qquad
\varrho(s)>0\quad(0<s<L).
\]
Let $S$ be the embedded surface obtained by revolving $\gamma$ about the $z$-axis.  If $L\le R$, then $D_R$ wraps the bounded solid enclosed by $S$ in Nandakumar's $1$-Lipschitz sense.
\end{lemma}

\begin{proof}
Set $q=L/R\le1$.  In polar coordinates $(r,\theta)$ on $D_R$, define
\[
F(r,\theta)=
\bigl(\varrho(qr)\cos\theta,\,\varrho(qr)\sin\theta,\,z(qr)\bigr).
\]
Since $\gamma$ is parametrized by arclength,
\[
\varrho'(s)^2+z'(s)^2=1
\]
for almost every $s$, so $|\varrho'|\le1$.  Because $\varrho(0)=0$,
\[
0\le\varrho(s)\le s.
\]
Away from the rotation axis the pullback metric is
\[
F^*g=q^2\,dr^2+\varrho(qr)^2\,d\theta^2
\le q^2\bigl(dr^2+r^2d\theta^2\bigr)
\le dr^2+r^2d\theta^2.
\]
Integrating lengths of curves and using continuity at the axis gives $\Lip(F)\le q\le1$.  The map is onto $S$, while the whole boundary circle $r=R$ is sent to the endpoint $\gamma(L)$ on the axis.  Hence the disk boundary is self-glued and the image is the closed surface $S$.
\end{proof}

\begin{proposition}[A soft answer to Question~1]\label[proposition]{prop:softQ1}
A disk of radius $\pi$ wraps a nonconvex solid of volume strictly greater than $4\pi/3$.
\end{proposition}

\begin{proof}
Consider the oblate spheroid with equatorial radius
\[
a=\frac{23}{20}
\]
and polar radius
\[
c=\frac45.
\]
Its volume is
\[
V_0=\frac{4\pi}{3}a^2c
=\frac{4\pi}{3}\frac{529}{500}
>\frac{4\pi}{3}.
\]
A north--south generating meridian is
\[
\gamma_0(t)=(a\sin t,c\cos t),\qquad 0\le t\le\pi.
\]
If $L_0$ denotes its length, Cauchy--Schwarz gives
\begin{align*}
L_0^2
&=\left(\int_0^\pi
\sqrt{a^2\cos^2t+c^2\sin^2t}\,dt\right)^2\\
&\le \pi\int_0^\pi
\bigl(a^2\cos^2t+c^2\sin^2t\bigr)\,dt\\
&=\pi^2\frac{a^2+c^2}{2}
=\pi^2\frac{157}{160}<\pi^2.
\end{align*}
Thus $L_0<\pi$ with strict slack.

We now make an arbitrarily small inward notch near the equator.  For $\delta>0$ small, put
\[
h=c\sin\delta,
\qquad
p_\pm=(a\cos\delta,\pm h).
\]
Replace the elliptical arc between $p_+$ and $p_-$ by the two line segments joining these points to
\[
v=(a\cos\delta-\varepsilon,0),
\qquad
0<\varepsilon<a\cos\delta.
\]
Because the ellipse is convex, the two new segments lie inside it.  Revolving the modified simple meridian gives an embedded solid $B_{\delta,\varepsilon}$ obtained from the spheroid by a rotational inward notch.

Its meridian length satisfies
\[
L_{\delta,\varepsilon}
\le L_0+2\sqrt{h^2+\varepsilon^2},
\]
so for sufficiently small $\delta$ and $\varepsilon$ we still have $L_{\delta,\varepsilon}<\pi$.  The removed volume lies inside the slab $|z|\le h$ and the cylinder $\varrho\le a$, hence
\[
0\le V_0-\Vol(B_{\delta,\varepsilon})
\le 2\pi a^2h.
\]
Since $h\to0$ with $\delta\to0$ and $V_0>4\pi/3$, the volume remains greater than $4\pi/3$ for sufficiently small $\delta$.

Finally, the meridional cross-section is nonconvex: it contains $p_+$ and $p_-$, but their midpoint $(a\cos\delta,0)$ lies outside the notched body because the new boundary radius at $z=0$ is $a\cos\delta-\varepsilon$.  Therefore $B_{\delta,\varepsilon}$ is nonconvex.  Applying \cref{lem:radial} with $R=\pi$ completes the proof.
\end{proof}

Thus the qualitative content of Question~1 follows without computer assistance.  The rest of the paper concerns the stronger quantitative problem of pushing the coefficient in \eqref{eq:Wdef} upward with a fully explicit certificate.

